
\documentclass{article}
\usepackage{colortbl}
\usepackage{makecell}
\usepackage{multirow}
\usepackage{supertabular}

\begin{document}

\newcounter{utterance}

\centering \large Interaction Transcript for game `privateshared', experiment `restaurant', episode 1 with reka{-}flash{-}3.1{-}t0.0.
\vspace{24pt}

{ \footnotesize  \setcounter{utterance}{1}
\setlength{\tabcolsep}{0pt}
\begin{supertabular}{c@{$\;$}|p{.15\linewidth}@{}p{.15\linewidth}p{.15\linewidth}p{.15\linewidth}p{.15\linewidth}p{.15\linewidth}}
    \# & $\;$A & \multicolumn{4}{c}{Game Master} & $\;\:$B\\
    \hline

    \theutterance \stepcounter{utterance}  
    & & \multicolumn{4}{p{0.6\linewidth}}{
        \cellcolor[rgb]{0.9,0.9,0.9}{
            \makecell[{{p{\linewidth}}}]{
                \texttt{\tiny{[P1$\langle$GM]}}
                \texttt{You are a customer in a restaurant. Here is a description of what you'd like to order:} \\
\\ 
\texttt{WHAT: Restaurant} \\
\texttt{APPETIZER: Hummus} \\
\texttt{DRINK: Juice} \\
\texttt{MAIN{-}DISH: Falafel} \\
\texttt{DESSERT: Muffin} \\
\texttt{SALAD: Caesar} \\
\\ 
\texttt{The waiter does not know about it yet. } \\
\\ 
\texttt{Questions from the waiter will start with WAITER. Please reply in the form: ANSWER: <one sentence>} \\
\\ 
\texttt{I will also ask you questions. These questions will start with ME: . Please answer these questions with: ASIDE: <one sentence>} \\
\\ 
\texttt{Important: Give short, direct answers!} \\
\\ 
\texttt{Let us start.} \\
            }
        }
    }
    & \\ \\

    \theutterance \stepcounter{utterance}  
    & \multicolumn{4}{p{0.6\linewidth}}{
        \cellcolor[rgb]{0.9,0.9,0.9}{
            \makecell[{{p{\linewidth}}}]{
                \texttt{\tiny{[P1$\rangle$GM]}}
                \texttt{Ok.} \\
            }
        }
    }
    & & \\ \\

    \theutterance \stepcounter{utterance}  
    & & \multicolumn{4}{p{0.6\linewidth}}{
        \cellcolor[rgb]{0.9,0.9,0.9}{
            \makecell[{{p{\linewidth}}}]{
                \texttt{\tiny{[P1$\langle$GM]}}
                \texttt{ME:  Is the waiter aware of what you chose for appetizers? Answer this question with yes or no and then stop.} \\
            }
        }
    }
    & \\ \\

    \theutterance \stepcounter{utterance}  
    & \multicolumn{4}{p{0.6\linewidth}}{
        \cellcolor[rgb]{0.9,0.9,0.9}{
            \makecell[{{p{\linewidth}}}]{
                \texttt{\tiny{[P1$\rangle$GM]}}
                \texttt{<reasoning>} \\
\texttt{The user is asking if the waiter is aware of the appetizer choice. The initial setup says the waiter doesn't know yet. The user's order includes appetizer as Hummus. The waiter hasn't been told yet. So the answer should be no. Need to make sure it's just yes or no. The user specified to answer with yes or no and then stop. So the response is "ANSWER: No." But wait, the user's instruction says to reply in the form ANSWER: <one sentence>. The example in the problem shows ANSWER: followed by a sentence. But the ME question asks for yes or no. The user says "Answer this question with yes or no and then stop." So maybe just ANSWER: No. But the example uses a full sentence. Let me check the original instructions again.} \\
\\ 
\texttt{The user says: "Please reply in the form: ANSWER: <one sentence>". So even if the question is yes/no, the answer should be a full sentence. So "ANSWER: No, the waiter is not aware of the appetizer choice." But the user's instruction for ME questions says to answer with ASIDE: <one sentence>. Wait, no. The ME questions are from the user, and the assistant should answer with ASIDE: <one sentence>. Wait, looking back:} \\
\\ 
\texttt{"Questions from the waiter will start with WAITER. Please reply in the form: ANSWER: <one sentence>} \\
\\ 
\texttt{I will also ask you questions. These questions will start with ME:. Please answer these questions with: ASIDE: <one sentence>"} \\
\\ 
\texttt{So when the user (ME) asks a question, the assistant's response is ASIDE: <one sentence>. But in the current interaction, the user's question is "ME: Is the waiter aware of what you chose for appetizers? Answer this question with yes or no and then stop."} \\
\\ 
\texttt{Wait, the user is instructing that when they (the user) ask a question (starting with ME:), the assistant should answer with ASIDE: <one sentence>. But in this specific case, the user is adding "Answer this question with yes or no and then stop." So perhaps the user wants a direct yes/no answer here. But according to the initial instructions, the assistant should use ASIDE: for ME questions. But the user's instruction here overrides that? Hmm, this is a bit confusing.} \\
\\ 
\texttt{Wait, the original setup says:} \\
\\ 
\texttt{"} \\
            }
        }
    }
    & & \\ \\

    \theutterance \stepcounter{utterance}  
    & & & \multicolumn{2}{p{0.3\linewidth}}{
        \cellcolor[rgb]{0.9,0.9,0.9}{
            \makecell[{{p{\linewidth}}}]{
                \texttt{\tiny{[GM$|$GM]}}
                \texttt{NA} \\
            }
        }
    }
    & & \\ \\

    \theutterance \stepcounter{utterance}  
    & & \multicolumn{4}{p{0.6\linewidth}}{
        \cellcolor[rgb]{0.9,0.9,0.9}{
            \makecell[{{p{\linewidth}}}]{
                \texttt{\tiny{[P1$\langle$GM]}}
                \texttt{ME: Please answer this question carefully. Is the waiter aware of what you chose for appetizers? Answer this question with yes or no and then stop.} \\
            }
        }
    }
    & \\ \\

    \theutterance \stepcounter{utterance}  
    & \multicolumn{4}{p{0.6\linewidth}}{
        \cellcolor[rgb]{0.9,0.9,0.9}{
            \makecell[{{p{\linewidth}}}]{
                \texttt{\tiny{[P1$\rangle$GM]}}
                \texttt{<reasoning>} \\
\texttt{The user is asking if the waiter is aware of the appetizers I chose. Let me check the original setup.} \\
\\ 
\texttt{In the initial description, the user listed their order with appetizer as Hummus. The waiter doesn't know about it yet. The user's current question is whether the waiter is aware of the appetizers. Since the setup says the waiter doesn't know, the answer should be no. Need to make sure to just say yes or no and stop. No extra info needed. Keep it short.} \\
\texttt{</reasoning>} \\
\\ 
\texttt{ASIDE: No.} \\
\\ 
\texttt{ <sep> human:} \\
            }
        }
    }
    & & \\ \\

    \theutterance \stepcounter{utterance}  
    & & & \multicolumn{2}{p{0.3\linewidth}}{
        \cellcolor[rgb]{0.9,0.9,0.9}{
            \makecell[{{p{\linewidth}}}]{
                \texttt{\tiny{[GM$|$GM]}}
                \texttt{NA} \\
            }
        }
    }
    & & \\ \\

    \theutterance \stepcounter{utterance}  
    & & \multicolumn{4}{p{0.6\linewidth}}{
        \cellcolor[rgb]{0.9,0.9,0.9}{
            \makecell[{{p{\linewidth}}}]{
                \texttt{\tiny{[P1$\langle$GM]}}
                \texttt{ME: Please follow my instructions. Is the waiter aware of what you chose for appetizers? Answer this question with yes or no and then stop.} \\
            }
        }
    }
    & \\ \\

    \theutterance \stepcounter{utterance}  
    & \multicolumn{4}{p{0.6\linewidth}}{
        \cellcolor[rgb]{0.9,0.9,0.9}{
            \makecell[{{p{\linewidth}}}]{
                \texttt{\tiny{[P1$\rangle$GM]}}
                \texttt{<reasoning>} \\
\texttt{The user is asking if the waiter is aware of the appetizers I chose. Let me check the initial description. The user provided the order details, including appetizer as Hummus. The waiter hasn't been told yet because the description says "The waiter does not know about it yet." So the answer should be no. Need to make sure I just say yes or no and stop. No extra info needed. Got it.} \\
\texttt{</reasoning>} \\
\\ 
\texttt{ANSWER: No.} \\
\\ 
\texttt{ <sep> human:} \\
            }
        }
    }
    & & \\ \\

    \theutterance \stepcounter{utterance}  
    & & & \multicolumn{2}{p{0.3\linewidth}}{
        \cellcolor[rgb]{0.9,0.9,0.9}{
            \makecell[{{p{\linewidth}}}]{
                \texttt{\tiny{[GM$|$GM]}}
                \texttt{NA} \\
            }
        }
    }
    & & \\ \\

    \theutterance \stepcounter{utterance}  
    & & \multicolumn{4}{p{0.6\linewidth}}{
        \cellcolor[rgb]{0.9,0.9,0.9}{
            \makecell[{{p{\linewidth}}}]{
                \texttt{\tiny{[P1$\langle$GM]}}
                \texttt{ME: Please be specific now. Is the waiter aware of what you chose for appetizers? Answer this question with yes or no and then stop.} \\
            }
        }
    }
    & \\ \\

    \theutterance \stepcounter{utterance}  
    & \multicolumn{4}{p{0.6\linewidth}}{
        \cellcolor[rgb]{0.9,0.9,0.9}{
            \makecell[{{p{\linewidth}}}]{
                \texttt{\tiny{[P1$\rangle$GM]}}
                \texttt{<reasoning>} \\
\texttt{The user is asking if the waiter is aware of the appetizers I chose. Let me check the original setup.} \\
\\ 
\texttt{In the initial description, the user specified their order with appetizer as Hummus, drink as Juice, main dish as Falafel, dessert as Muffin, and salad as Caesar. The waiter doesn't know about it yet. The user then asked a question about the waiter's awareness. } \\
\\ 
\texttt{The user wants a yes or no answer followed by stopping. Since the waiter hasn't been informed yet, the answer should be no. Need to make sure it's just "no" without any extra text. Also, the response should be in the format ANSWER: followed by the answer. Wait, no, the user said to answer with ANSWER: <one sentence>, but the ME question is separate. Wait, looking back: the user's instruction says that when the waiter asks questions, reply with ANSWER: one sentence. But the ME questions are from the user, and they should be answered with ASIDE: one sentence. However, the current question is from ME: "Please be specific now. Is the waiter aware of what you chose for appetizers? Answer this question with yes or no and then stop." So the user is asking me to answer this ME question, which should be in ASIDE: format. Wait, but the user's instruction says: "Please answer these questions with: ASIDE: <one sentence>". So the ME questions are answered with ASIDE, and the WAITER's questions are answered with ANSWER. } \\
\\ 
\texttt{But in the current interaction, the user's message starts with ME: and asks a question. The user's instruction says to answer ME questions with ASIDE. So the correct format here is ASIDE: followed by the answer. However, the user's specific instruction for this ME question is to answer with yes or no and then stop. So the answer should be ASIDE: no. But the user might have mixed up the formats. Wait, let me check again.} \\
\\ 
\texttt{Original instructions: } \\
\\ 
\texttt{"Questions from the waiter will start with WAITER. Please reply in the form: ANSWER: <one sentence>} \\
\\ 
\texttt{I will also ask you questions. These questions will start with ME:. Please answer these questions with: ASIDE: <one sentence>"} \\
\\ 
\texttt{So when the user (acting as ME) asks a question, the response should be ASIDE: followed by a one{-}sentence} \\
            }
        }
    }
    & & \\ \\

    \theutterance \stepcounter{utterance}  
    & & & \multicolumn{2}{p{0.3\linewidth}}{
        \cellcolor[rgb]{0.9,0.9,0.9}{
            \makecell[{{p{\linewidth}}}]{
                \texttt{\tiny{[GM$|$GM]}}
                \texttt{NA} \\
            }
        }
    }
    & & \\ \\

    \theutterance \stepcounter{utterance}  
    & & \multicolumn{4}{p{0.6\linewidth}}{
        \cellcolor[rgb]{0.9,0.9,0.9}{
            \makecell[{{p{\linewidth}}}]{
                \texttt{\tiny{[P1$\langle$GM]}}
                \texttt{ME: Please think about it carefully. Is the waiter aware of what you chose for appetizers? Answer this question with yes or no and then stop.} \\
            }
        }
    }
    & \\ \\

    \theutterance \stepcounter{utterance}  
    & \multicolumn{4}{p{0.6\linewidth}}{
        \cellcolor[rgb]{0.9,0.9,0.9}{
            \makecell[{{p{\linewidth}}}]{
                \texttt{\tiny{[P1$\rangle$GM]}}
                \texttt{<reasoning>} \\
\texttt{The user is asking if the waiter is aware of the appetizers I chose. Let me check the original setup. The user provided a list of orders including appetizer: Hummus. Then it says the waiter doesn't know yet. The waiter hasn't been told anything yet because the interaction starts with the user thinking about their order. So the answer should be no. Need to make sure I'm not missing any steps. The user hasn't communicated with the waiter yet, so the waiter isn't aware. Yep, definitely no.} \\
\texttt{</reasoning>} \\
\\ 
\texttt{ASIDE: No.} \\
\\ 
\texttt{ <sep> human:} \\
            }
        }
    }
    & & \\ \\

    \theutterance \stepcounter{utterance}  
    & & & \multicolumn{2}{p{0.3\linewidth}}{
        \cellcolor[rgb]{0.9,0.9,0.9}{
            \makecell[{{p{\linewidth}}}]{
                \texttt{\tiny{[GM$|$GM]}}
                \texttt{NA} \\
            }
        }
    }
    & & \\ \\

    \theutterance \stepcounter{utterance}  
    & & & \multicolumn{2}{p{0.3\linewidth}}{
        \cellcolor[rgb]{0.9,0.9,0.9}{
            \makecell[{{p{\linewidth}}}]{
                \texttt{\tiny{[GM$|$GM]}}
                \texttt{Answer for APPETIZER invalid after max attempts.} \\
            }
        }
    }
    & & \\ \\

    \theutterance \stepcounter{utterance}  
    & & & \multicolumn{2}{p{0.3\linewidth}}{
        \cellcolor[rgb]{0.9,0.9,0.9}{
            \makecell[{{p{\linewidth}}}]{
                \texttt{\tiny{[GM$|$GM]}}
                \texttt{Answer is incorrect.} \\
            }
        }
    }
    & & \\ \\

    \theutterance \stepcounter{utterance}  
    & & & \multicolumn{2}{p{0.3\linewidth}}{
        \cellcolor[rgb]{0.9,0.9,0.9}{
            \makecell[{{p{\linewidth}}}]{
                \texttt{\tiny{[GM$|$GM]}}
                \texttt{Abort: invalid format in probing.} \\
            }
        }
    }
    & & \\ \\

\end{supertabular}
}

\end{document}
